% !TEX root = _yoko.tex

% ==========================================
% プロジェクト予稿（1P）の本文
% ==========================================

\section{はじめに}
音楽や物語といった作品を分析する際には，
出来事を分解し，それらの関係性を整理しながら，全体構造を捉えることが重要となる．
このような分析では，出来事同士のまとまりや重要度を検討し，
それらを階層構造として表現する作業がしばしば用いられる．

このような構造化を中心とする分析手法の一つが，
EMDA（Emotion Movement Design Annotator）分析である\cite{emda2019rigolo}．
EMDA分析では，音楽や物語を構成する出来事を「イベント」として抽出し，
それらのまとまりと重要度を判断しながら，
作品の根幹に接続する形で木構造として整理する．
この木構造は，単なる図示ではなく，
分析者の解釈判断そのものを可視化する表現である．

一方で，EMDA分析においては，
解釈の更新に伴って木構造を何度も更新する必要があり，
その都度，線や配置を手作業で修正することが大きな負担となる．

そこで本研究では，
EMDA分析の工程のうち，判断結果を木構造として外化する作図工程に着目し，
構造作成・修正を低コストで反復できる編集支援ツールの開発を行った．
本ツールは，分析結果を自動生成したり，
判断の正誤を評価したりするものではない．
あくまで，分析者が自ら下した判断を前提として，
それを素早く組み立て直せる編集環境を提供することを目的としている．

\section{ツール概要}
本ツールでは，分析対象の各要素を線分として表現し，
線分同士のスナップ接続によって親子関係を構築するUIを採用した．
線分の上端をドラッグ操作で移動させることで，
従属関係の付与・解除を直感的に行うことができる．

親子関係は画面上の絶対位置ではなく，
親線分上の相対位置として保持されるため，
表示範囲の変更や再配置が行われても意味的関係が失われない．
また，親線の移動に子線が追従するため，
従属関係を変更する際にも全体を作り直す必要がない．

さらに，線分間の区間にイベント記述用テキストを配置し，
その下段に複数イベントをまとめるグループ単位の記述欄を設けた．
これにより，個別イベントの記述と，
それらを束ねた解釈を異なる粒度で保持できる．

加えて，線の追加・削除，
グループ境界の分割・結合，
Undo/Redo，状態保存といった機能を実装し，
構造の試行錯誤を前提とした分析作業を支援している．

\figref{fig:01}に，
本ツールを用いて多数のイベントを含む木構造を作成した例を示す．
構造と記述が一体となって配置されていることが分かる．

\begin{figure}[t]
  \centering
  \includegraphics[width=0.95\linewidth]{figures/fig01.png}
  \caption{本ツールで作成した木構造の例}\label{fig:01}
\end{figure}



\section{操作性の評価}
評価として，PowerPointと本ツールを用いて
ランダム生成した簡易的な木構造を作成する定量比較を行った．
その結果，本ツールを用いた場合，
木構造の作成に要する時間が概ね2～3倍程度短縮される傾向が観察された．

また，同様の構造を持つ木構造に対して，
従属関係の変更やイベントの追加・削除といった修正作業を行い，
操作過程を定性的に比較した．
その結果，PowerPointでは線一本ずつの再配置や全体の調整が必要となる場面でも，
本ツールでは親子関係に基づく追従機構により，
最小限の操作で構造更新が可能であることが観察された．

これらの結果は，
本ツールが単に作図時間を短縮するだけでなく，
構造修正に伴う手続き的負担を軽減し，
分析者が構造の試行錯誤を行いやすい環境を提供していることを示している．

\section{おわりに}
本研究では，EMDA分析において特に負担となりやすい「木構造の作成」に着目し，
その簡略化および作業時間の短縮を目的とした支援ツールを設計・実装した．

評価の結果，従来PowerPoint等で手作業により行っていた構造作成および修正作業を，
比較的少ない操作で反復的に行えることを確認した．
一方で，本ツールは分析作業全体を完結させる完成形の環境を提供するものではなく，
構造操作の負担軽減に焦点を当てた初期的・試作的な実装に留まっている．
今後は，評価結果や利用状況を踏まえ，操作性や機能の改善を通じて，
分析過程により適合した支援環境へと発展させていく必要がある．

本ツールが，分析者の作業効率を高めるだけでなく，
これまでEMDA分析に取り組む機会の少なかった利用者にとっても，
分析への導入障壁を下げる一助となることを期待する．